% libre-compendium
%
% Copyright (C) 2026 Richard Qin
%
% This document is dual-licensed under either:
% - GNU Free Documentation License, Version 1.3 or any later version
%   published by the Free Software Foundation; with no Invariant Sections,
%   no Front-Cover Texts, and no Back-Cover Texts (see COPYING.DOC); or
% - Creative Commons Attribution-ShareAlike 4.0 International License
%   (see LICENSE-CC-BY-SA-4.0).

% Extracted from chapter03-mechanical-wave.tex.
\begin{exercise}[level=basic]
说明机械波产生的两个条件，并说明机械波传播时介质质元是否随波迁移。
\begin{solution}
需要波源和弹性介质。机械波传播的是振动状态、相位和能量，介质质元只在各自平衡位置附近振动，不随波作永久迁移。
\end{solution}
\end{exercise}

\begin{exercise}[level=basic]
区分横波和纵波，并说明空气中的声波属于哪一类。
\begin{solution}
横波中质元振动方向垂直于传播方向；纵波中质元振动方向平行于传播方向。空气中的声波为纵波。
\end{solution}
\end{exercise}

\begin{exercise}[level=basic]
从是否依赖弹性介质、扰动是否持续、传播区域三个角度，各举出一种波的分类。
\begin{solution}
示例：按是否依赖弹性介质，可分为机械波和非机械波；按扰动是否持续，可分为脉冲波和连续波；按传播区域，可分为体波和表面波。
\end{solution}
\end{exercise}

\begin{exercise}[level=basic]
写出沿 $x$ 正向传播的平面简谐波函数，并说明 $A,\omega,k,\lambda,T,\nu,u$ 的含义和关系。
\begin{solution}
沿 $x$ 正向传播：
\[
  y=A\cos(\omega t-kx+\varphi_0).
\]
$A$ 为振幅，$\omega=2\pi/T=2\pi\nu$，$k=2\pi/\lambda$，
\[
  u=\lambda\nu=\frac{\omega}{k}.
\]
\end{solution}
\end{exercise}

\begin{exercise}[level=basic]
已知 $y=A\cos(\omega t-kx+\varphi_0)$，证明它沿 $x$ 正向传播。
\begin{solution}
取相位为常数：
\[
  \omega t-kx+\varphi_0=C.
\]
解得
\[
  x=\frac{\omega}{k}t+\frac{\varphi_0-C}{k}.
\]
同一相位点随时间向 $x$ 增大的方向移动，故波沿 $x$ 正向传播。
\end{solution}
\end{exercise}

\begin{exercise}[level=standard]
一列波 $\lambda=0.50\,\mathrm{m}$、$\nu=20\,\mathrm{Hz}$，求波速、周期、角频率和角波数。
\begin{solution}
\[
  u=\lambda\nu=0.50\times20=10\,\mathrm{m/s},
\]
\[
  T=\frac1\nu=0.05\,\mathrm{s},\qquad
  \omega=2\pi\nu=40\pi\,\mathrm{rad/s},
\]
\[
  k=\frac{2\pi}{\lambda}=4\pi\,\mathrm{rad/m}.
\]
\end{solution}
\end{exercise}

\begin{exercise}[level=standard]
对沿 $x$ 正向传播的平面简谐波，求同一时刻相距 $\lambda/3$ 的两点相位差大小，并说明沿传播方向较远点相位超前还是落后。
\begin{solution}
相位差大小为
\[
  |\Delta\varphi|=\frac{2\pi}{\lambda}\frac{\lambda}{3}
  =\frac{2\pi}{3}.
\]
沿 $x$ 正向传播时，传播方向上更远的点相位落后。
\end{solution}
\end{exercise}

\begin{exercise}[level=standard]
原点振动方程为 $y=0.02\cos(4\pi t+\pi/3)\,\mathrm{m}$，波沿 $x$ 正向传播，波速 $u=2\,\mathrm{m/s}$。写出波函数。
\begin{solution}
\[
  y(x,t)=0.02\cos\left[4\pi\left(t-\frac{x}{2}\right)+\frac{\pi}{3}\right]
  =0.02\cos\left(4\pi t-2\pi x+\frac{\pi}{3}\right)\,\mathrm{m}.
\]
\end{solution}
\end{exercise}

\begin{exercise}[level=standard]
某波函数为 $y=0.05\cos(10\pi t+2\pi x)\,\mathrm{m}$，判断传播方向、波长、频率和波速。
\begin{solution}
因相位为 $10\pi t+2\pi x$，波沿 $x$ 负向传播。
\[
  \omega=10\pi,\qquad k=2\pi.
\]
故
\[
  \nu=\frac{\omega}{2\pi}=5\,\mathrm{Hz},\qquad
  \lambda=\frac{2\pi}{k}=1\,\mathrm{m},
\]
\[
  u=\lambda\nu=5\,\mathrm{m/s}.
\]
\end{solution}
\end{exercise}

\begin{exercise}[level=standard]
沿 $x$ 正向传播的横波在某点波形斜率为正，判断该质元瞬时速度方向。
\begin{solution}
沿 $x$ 正向传播时
\[
  \frac{\partial y}{\partial t}=-u\frac{\partial y}{\partial x}.
\]
斜率为正，则 $\partial y/\partial t<0$，质元向负 $y$ 方向运动。
\end{solution}
\end{exercise}

\begin{exercise}[level=standard]
写出一维波动方程，并说明波函数和波动方程的区别。
\begin{solution}
一维波动方程：
\[
  \frac{\partial^2 \xi}{\partial x^2}
  =\frac1{u^2}\frac{\partial^2 \xi}{\partial t^2}.
\]
波动方程是介质的动力学方程；波函数是满足该方程并由初始、边界条件决定的具体解。
\end{solution}
\end{exercise}

\begin{exercise}[level=standard]
两同相相干波源到某点路程差为 $3\lambda$，判断干涉加强或减弱。
\begin{solution}
路程差为整数倍波长，同相波源在该点相位差为 $2m\pi$，干涉加强。
\end{solution}
\end{exercise}

\begin{exercise}[level=standard]
两同相相干波源到某点路程差为 $2.5\lambda$，判断干涉加强或减弱。
\begin{solution}
路程差为半整数倍波长，相位差为奇数倍 $\pi$，干涉减弱。
\end{solution}
\end{exercise}

\begin{exercise}[level=standard]
两相干波在点 $P$ 的振幅分别为 $A_1,A_2$，相位差为 $\Delta\varphi$，写出合振幅公式。
\begin{solution}
\[
  A=\sqrt{A_1^2+A_2^2+2A_1A_2\cos\Delta\varphi}.
\]
\end{solution}
\end{exercise}

\begin{exercise}[level=standard]
两同相波源频率 $50\,\mathrm{Hz}$，波速 $10\,\mathrm{m/s}$。点 $P$ 到两波源距离差 $0.30\,\mathrm{m}$，判断干涉情况。
\begin{solution}
\[
  \lambda=\frac{u}{\nu}=\frac{10}{50}=0.20\,\mathrm{m}.
\]
路程差 $0.30\,\mathrm{m}=1.5\lambda$，故干涉减弱。
\end{solution}
\end{exercise}

\begin{exercise}[level=standard]
声源频率 $800\,\mathrm{Hz}$，声速 $340\,\mathrm{m/s}$。声源静止，观察者以 $10\,\mathrm{m/s}$ 向声源运动，求接收频率。
\begin{solution}
声源静止，观察者向声源运动：
\[
  \nu_R=\nu_s\frac{u+v_R}{u}
  =800\frac{340+10}{340}\approx823.5\,\mathrm{Hz}.
\]
\end{solution}
\end{exercise}

\begin{exercise}[level=standard]
声源频率 $800\,\mathrm{Hz}$，观察者静止，声源以 $20\,\mathrm{m/s}$ 远离观察者，求接收频率。
\begin{solution}
观察者静止，声源远离，取 $v_S=-20\,\mathrm{m/s}$：
\[
  \nu_R=800\frac{340}{340+20}\approx755.6\,\mathrm{Hz}.
\]
\end{solution}
\end{exercise}

\begin{exercise}[level=advanced]
波沿 $x$ 正向传播，$t=0$ 时原点位移为零且向正方向运动。振幅 $A$、角频率 $\omega$、波数 $k$ 已知，写出波函数。
\begin{solution}
设原点振动为 $y(0,t)=A\cos(\omega t+\varphi_0)$。$t=0$ 时 $y=0$ 且速度为正：
\[
  \cos\varphi_0=0,\qquad -A\omega\sin\varphi_0>0.
\]
故可取 $\varphi_0=-\pi/2$。波沿 $x$ 正向传播：
\[
  y=A\cos\left(\omega t-kx-\frac{\pi}{2}\right).
\]
\end{solution}
\end{exercise}

\begin{exercise}[level=advanced]
两相干波源初相差为 $\pi$，到某点路程差为 $m\lambda$。判断该点加强还是减弱。
\begin{solution}
路程差为 $m\lambda$ 本会给出 $2m\pi$ 相位差；但波源初相差为 $\pi$，总相位差为奇数倍 $\pi$，故减弱。
\end{solution}
\end{exercise}

\begin{exercise}[level=standard]
写出驻波 $y=2A\cos kx\cos\omega t$ 的波节、波腹位置，并给出相邻波节间距。
\begin{solution}
波节满足
\[
  \cos kx=0,
\]
即
\[
  kx=\left(n+\frac12\right)\pi,\qquad
  x=\left(n+\frac12\right)\frac{\lambda}{2}
  =\frac{(2n+1)\lambda}{4}.
\]
波腹满足
\[
  |\cos kx|=1,
\]
即
\[
  x=\frac{n\lambda}{2}.
\]
相邻波节间距为 $\lambda/2$，相邻波节与波腹间距为 $\lambda/4$。
\end{solution}
\end{exercise}

